\chapter{La crisis}
\label{ch:crisis}

\section{Una nota sobre escribir del presente}

Este capítulo es el que más se acerca al día en que se escribió, y cuanto más se
acerca menos certeza tiene. La publicación estadística cubana se adelgazó después
de 2020; varias series del \emph{Anuario} dejaron de aparecer o aparecieron con
mucho retraso; y las cifras de 2025 y 2026 que circulan son en la mayoría de los
casos estimaciones de economistas y demógrafos independientes y no datos
publicados. Donde este capítulo da una cifra reciente dice de quién es la cifra, y
los números de los dos últimos años deben leerse como provisionales y contrastarse
con lo que se haya publicado desde entonces. La dirección de todos ellos, en
cambio, no está en duda.

\section{Un presidente nuevo y una constitución nueva}

Miguel Díaz-Canel asumió la presidencia en abril de 2018 y la Primera Secretaría
del partido en abril de 2021, cuando Raúl Castro se retiró. Es el primer dirigente
cubano desde 1959 que no se apellida Castro y el primero que no combatió en la
insurrección, y ninguno de los dos hechos se ha traducido en mucha autoridad
propia.

Su primer acto sustantivo fue una constitución. El texto de 2019, aprobado en
referendo en febrero, es un documento más liberal que el de 1976 al que sustituyó,
y esto conviene decirlo antes que el matiz. Reconoce la propiedad privada y la
inversión extranjera como componentes de la economía y no como excepciones
toleradas; restablece la presunción de inocencia y el habeas corpus; limita al
presidente a dos mandatos de cinco años e impone un límite de edad para la primera
elección; recrea el cargo de primer ministro; refuerza el lenguaje sobre no
discriminación. En contra, conserva al Partido Comunista como fuerza dirigente
superior de la sociedad y del Estado, y conserva la irrevocabilidad del socialismo
que se había insertado en 2002 para enterrar el Proyecto Varela. Es una reforma de
la administración del Estado y no de su política.

El borrador había contenido una definición de matrimonio que habría permitido el
matrimonio entre personas del mismo sexo; se retiró tras la oposición organizada de
las iglesias evangélicas y se remitió a un futuro código de familia. Ese código se
sometió a referendo aparte en septiembre de 2022 y se aprobó con cerca de dos
tercios de los votos, lo que hace de Cuba uno de los pocos países del mundo que ha
legalizado el matrimonio igualitario por voto popular. Es un avance liberal genuino
y llegó por referendo, lo que vale la pena señalar en un capítulo por lo demás sin
alivio.

\section{El reapriete}

Entre 2017 y 2020 Estados Unidos impuso una larga serie de medidas ---la cifra que
suele citarse es 243--- que revirtieron la apertura y fueron más allá.\unv{} Las
escalas de cruceros se prohibieron en junio de 2019. Los viajes individuales de
pueblo a pueblo terminaron. Las remesas se limitaron y después, a finales de 2020,
quedaron efectivamente cortadas cuando Western Union cerró sus más de cuatrocientas
oficinas cubanas tras la designación por el Tesoro norteamericano de la procesadora
vinculada a los militares por la que corrían. El Título III de la Helms-Burton
entró en vigor en mayo de 2019 por primera vez desde 1996, exponiendo a cualquier
empresa que usara propiedad confiscada a ser demandada en tribunales
norteamericanos. Y en enero de 2021, en la última semana de la administración
Trump, Cuba volvió a la lista de Estados patrocinadores del terrorismo,
designación cuyo efecto práctico está menos en sus propias prohibiciones que en la
conducta que induce en los bancos, que se protegen negándose por completo a
cualquier negocio cubano.

La administración Biden alivió partes de esto desde 2022 ---topes a las remesas
levantados, viajes en grupo restaurados, procesamiento consular reanudado en La
Habana--- y en enero de 2025 retiró la designación de terrorismo, que se
restableció a los pocos días del cambio de administración. El efecto neto de la
década es que el entorno externo se endureció, que el endurecimiento se concentró
exactamente en los flujos que llegaban a los hogares cubanos y no al Estado cubano
---remesas, viajes independientes, alojamiento privado--- y que la política cubana
hacia su propio sector privado, descrita en el capítulo anterior, había empezado a
endurecerse primero.

Entretanto el subsidio venezolano se encogió con Venezuela. Las entregas de
petróleo, que habían rondado los 100.000 barriles diarios, cayeron a una fracción
de eso, cubriéndose el faltante de forma irregular con Rusia, México y compras
puntuales que el país no podía pagar con fiabilidad.

\section{Dos años con las fronteras cerradas}

El turismo había llegado a unos 4,7 millones de visitantes en 2018. Cayó a unos
1,1 millones en 2020 y a algo menos de 400.000 en 2021.\unv{} Para una economía en
la que el turismo era la segunda o tercera fuente de divisas y la primera fuente de
demanda del sector privado, cerrar las fronteras dos años equivalía a cerrar la
economía privada. Las llegadas se recuperaron hasta unos 2,4 millones para 2023 y
después se estancaron por debajo de esa cifra, en torno a la mitad del nivel
prepandémico, mientras el Caribe en conjunto se recuperaba del todo ---divergencia
que tiene más que ver con los apagones, los hoteles sin agua y las tiendas vacías
que con la política norteamericana---.\unv{}

Frente a esto, la pandemia produjo el único éxito cubano inequívoco del período. El
sector biotecnológico cubano desarrolló sus propias vacunas contra la COVID-19
---Abdala en el Centro de Ingeniería Genética y Biotecnología, y Soberana~02 y
Soberana Plus en el Instituto Finlay---, hizo sus propios ensayos, reportó
eficacias por encima del 90 por ciento, fabricó a escala nacional y vacunó a más
del 90 por ciento de la población con producto nacional, incluidos los niños desde
los dos años antes que en la mayor parte del mundo.\unv{} Ningún otro país de
América Latina hizo esto. Cuba lo hizo sin poder importar jeringuillas de forma
fiable.

El matiz va con ello: las vacunas no se presentaron, o no completaron, las vías
internacionales de precalificación y regulación que habrían permitido venderlas a
escala en el exterior, que es el mismo fallo que el capítulo~\ref{ch:life}
identifica en todo el sector ---ciencia excelente, ninguna conversión regulatoria,
y por tanto muy poco dinero---. Cuba se vacunó a sí misma y no pudo vender el
resultado.

\section{La Tarea Ordenamiento}

El 1 de enero de 2021 ---en plena pandemia, con las fronteras cerradas, el turismo
en cero y la economía habiéndose contraído cerca del 11 por ciento el año
anterior--- el gobierno ejecutó la unificación monetaria que venía preparando desde
2011.

Se abolió el peso convertible con una ventana de conversión de seis meses. La tasa
oficial para las empresas pasó de un peso por dólar a veinticuatro, una devaluación
del 2.300 por ciento aplicada de golpe a todos los balances del país. Los salarios
y las pensiones se multiplicaron varias veces para compensar, con el salario mínimo
llevado a 2.100 pesos. Se recortaron subsidios y se subieron muchos precios
regulados.

La dirección era correcta y llevaba veinte años siéndolo. El
capítulo~\ref{ch:dollarreform} explicó por qué: con dos tasas separadas por un
factor veinticinco, ninguna contabilidad empresarial cubana significaba nada, y
nada podía reformarse porque nada podía medirse. La ejecución fue catastrófica, y
las razones son lo bastante específicas como para ser instructivas.

Se hizo sin reservas. Una tasa fija es la promesa de vender dólares a esa tasa, y
Cuba no tenía dólares, así que la promesa no se cumplió y apareció de inmediato una
tasa paralela.

Se hizo sin respuesta de la oferta. Se multiplicaron los salarios en una economía
donde la cantidad de bienes estaba cayendo, que es la definición de una inflación.

Se hizo con el sector privado todavía reprimido ---la congelación de licencias de
2017 seguía vigente y la ley de pequeñas y medianas empresas estaba a ocho meses de
distancia---, de modo que a la única parte de la economía que podía haber respondido
a precios más altos con más producción no se le permitió hacerlo.

Y se hizo en el fondo del peor choque externo desde 1993, cuando la secuencia
correcta es estabilizar primero y unificar después.

Los resultados se siguieron. La inflación oficial de precios al consumidor fue de
alrededor del 77 por ciento en 2021 y se mantuvo en dos dígitos altos los dos años
siguientes, con estimaciones independientes considerablemente mayores para los
bienes que la gente compra de verdad.\unv{} La tasa de cambio informal, que la
reforma pretendía eliminar, se despegó de la oficial y siguió subiendo: pasó de 100
pesos por dólar en 2022, de 200 en 2023, de 300 en 2024, y más desde
entonces.\unv{} Los salarios y pensiones reales terminaron por debajo de donde
estaban antes de la reforma que los había multiplicado. Los ahorros de los cubanos
se destruyeron por segunda vez en treinta años.

El capítulo~\ref{ch:money} está escrito contra este capítulo. La lección no es que
la unificación estuviera mal. Es que una reforma monetaria sin reservas, sin
respuesta de la oferta y sin un programa de estabilización por delante no es una
reforma sino una devaluación de los ahorros de la población, y que el mismo Estado
había demostrado en 1994--96, con menos recursos y mayor competencia, cómo se hace
en el orden correcto.

\section{11 de julio de 2021}

La mañana del domingo 11 de julio de 2021 la gente salió a la calle en San Antonio
de los Baños, un pueblo de cuarenta mil habitantes al suroeste de La Habana. Se
filmó y se subió a Facebook. En horas había manifestaciones en decenas de pueblos y
ciudades de toda la isla ---los conteos van de unas cuarenta a más de sesenta
localidades--- incluidas La Habana, Santiago, Camagüey, Holguín, Cienfuegos,
Matanzas y Palma Soriano.\unv{}

Nada semejante había ocurrido desde 1959. El maleconazo de 1994 fue una ciudad en
una tarde. Esto fue nacional, simultáneo y organizado por nadie ---que es
exactamente lo que lo hizo incontenible y, para el gobierno, incomprensible---.
Las causas inmediatas fueron los apagones en una ola de calor de julio, las
farmacias vacías durante la peor ola de COVID que tuvo Cuba, las colas de comida y
el choque de precios del ordenamiento. Los manifestantes eran abrumadoramente
jóvenes, pobres y desproporcionadamente de los barrios negros y mestizos que el
capítulo~\ref{ch:dollarreform} identificó como los hogares sin remesas ---es
decir, del sector social que la revolución siempre ha contado como suyo---.

La respuesta fue el corte nacional del internet móvil en cuestión de horas; la
orden de combate televisada del presidente, llamando a los revolucionarios a la
calle; el despliegue de fuerzas especiales y de brigadas civiles; y las
detenciones masivas. Un manifestante, Diubis Laurencio Tejeda, murió por disparos
en La Güinera, La Habana. Cubalex y Justicia~11J \citep{justicia11j}, los dos
grupos de observación que hicieron el conteo porque nadie oficial lo hizo,
documentaron 1.481 personas privadas de libertad en el contexto de las protestas,
de las cuales 57 eran menores de dieciocho años. Unas 600 fueron procesadas; al
menos 141 fueron acusadas de sedición; y las condenas fueron de cuatro a treinta
años. Tres años después, varios cientos seguían presos. Los procesos no fueron
contra una organización, porque no había organización; fueron contra individuos
que habían caminado por la calle.

Este es el hecho político central del período y la Parte III tiene que responder
ante él. Un gobierno cuya pretensión de legitimidad descansa en haber sacado a los
pobres y a los negros de la desatención de la república respondió a una
manifestación de los pobres y los negros con condenas de treinta años. Sea lo que
sea que el diseño de la Parte III proponga sobre reconciliación, amnistía y
secuencia de la apertura política, el 11 de julio de 2021 y sus presos son el caso
que tiene que poder resolver.

\section{La retirada del propio Estado}

En agosto de 2021, seis semanas después de las protestas, el gobierno hizo lo que
llevaba una década negándose a hacer: los Decretos-Leyes 46 a 49 crearon las
pequeñas y medianas empresas con personalidad jurídica. Un negocio privado cubano
podía por fin ser una empresa ---poseer activos, firmar contratos, contratar, e
importar y exportar a través de un intermediario estatal---. Se aprobaron unas once
o doce mil en los tres años siguientes, la gran mayoría privadas, y rápidamente
pasaron a ser una parte significativa de las importaciones de alimentos y bienes de
consumo del país.\unv{}

La reforma fue real y fue también, con toda transparencia, una rendición: el Estado
legalizó a los importadores privados porque ya no podía surtir sus propias tiendas.
La misma lógica produjo la redolarización parcial del comercio minorista a partir de
2024 ---nuevas tiendas estatales que venden solo en dólares, para captar remesas---
lo que reintrodujo, tres años después de abolirla, una economía minorista de dos
monedas con la moneda de los pobres del lado equivocado.

Y las empresas llegaron a un entorno de política que se volvió contra ellas casi de
inmediato. Desde 2024 llegaron los topes de precios a los comerciantes privados,
las restricciones a lo que podía importarse, impuestos más altos, y un argumento
público de la dirigencia según el cual las nuevas empresas eran la causa de la
inflación y no una respuesta a ella. El patrón de 1996 y de 2017 se repitió en un
ciclo de dos años.

\section{Las luces}

El sistema eléctrico cubano llegó al final de su vida en 2024.

El parque generador es un número pequeño de grandes unidades térmicas de fueloil,
construidas en su mayoría entre los años sesenta y ochenta, funcionando más allá de
su vida de diseño con crudo nacional pesado para el que no fueron construidas, sin
mantenimiento, sin repuestos y con frecuencia sin combustible. A lo largo de 2023 y
2024 los déficits diarios de un tercio o más de la demanda eran rutinarios, y los
apagones de doce a veinte horas eran normales fuera de La Habana y cada vez más
normales dentro.

La escala de la retirada se ve en la producción. Cuba generó del orden de 20
teravatios-hora en 2018; en 2023 generó 15,3, de los cuales 3,7 ---cerca de una
cuarta parte--- se perdieron en transmisión y distribución antes de llegar a
nadie.\unv{} El consumo por habitante cayó de unos 1.856 kilovatios-hora en 2018 a
1.387 en 2023.

El 18 de octubre de 2024 la central Antonio Guiteras, la mayor unidad individual del
país, se disparó, y toda la red nacional colapsó. Cuba quedó a oscuras: once
millones de personas, todas las provincias, sin excepciones. El restablecimiento
tomó varios días y falló dos veces durante el intento. Volvió a ocurrir en
noviembre tras el huracán Rafael, y otra vez en diciembre, y otra vez en
2025.\unv{} Un apagón nacional no es una escasez; es el fallo del sistema en cuanto
sistema, y una vez que una red ha colapsado colapsa con más facilidad.

La respuesta del gobierno es la energía solar, financiada y construida por China: un
programa anunciado hacia 2024 de decenas de parques fotovoltaicos con una primera
meta cercana a los mil megavatios, de los cuales se informó que varios cientos de
megavatios se energizaron durante 2025.\unv{} La dirección es correcta ---el
capítulo~\ref{ch:energy} sostiene que es la única dirección disponible--- pero la
capacidad fotovoltaica sin almacenamiento no atiende un pico de demanda que ocurre
después de la puesta del sol, y no vuelve a poner en marcha por sí sola un parque
térmico que sigue teniendo que cubrir la noche.

Todo lo demás en este capítulo está aguas abajo de las luces. Sin electricidad
estable no hay bombeo de agua, ni refrigeración, ni procesamiento de alimentos, ni
turismo digno del nombre, ni hospital que pueda planificar un salón, ni trabajo
remoto, ni industria. El capítulo~\ref{ch:premises} trata por eso la electricidad
como la restricción vinculante de todo el diseño y la pone por delante de cualquier
otro sector.

\section{El éxodo}

La mayor emigración de la historia cubana ocurrió entre 2021 y 2025, y empequeñece
a las que este libro ya ha descrito.

La ruta pasaba por Nicaragua, que eliminó el requisito de visa para los cubanos en
noviembre de 2021, y de ahí por tierra hasta la frontera de Estados Unidos. En el
año fiscal norteamericano 2022 las autoridades registraron algo más de 220.000
encuentros con nacionales cubanos, y unos 200.000 al año siguiente; un programa de
parole humanitario abierto en enero de 2023 admitió por aire a decenas de miles
más; y un número sustancial se fue a España al amparo de una ley de nacionalidad
por descendencia de 2022, y a México, Uruguay, Serbia y otros destinos.\unv{} Las
estimaciones razonables del total de salidas entre 2021 y 2025 van de alrededor de
un millón a cerca de dos, frente a una población que había sido de unos 11,2
millones.

La oficina de estadísticas oficial situó la población residente en 9,7 millones en
2024 ---una caída de 1,4 millones en cuatro años, y la primera vez por debajo de
diez millones desde los años ochenta--- con el censo de 2024 aplazado. Los
demógrafos cubanos independientes, trabajando con datos de nacimientos, defunciones
y fronteras, sitúan la cifra real de residentes aún más abajo, un estudio en torno
a ocho millones, y Albizu-Campos estima unas 1,79 millones de salidas solo entre
2021 y 2024 \citep{onei}.\unv{} La incertidumbre es grande y el hecho no lo es:
Cuba ha perdido algo así como entre una décima y una quinta parte de su gente en
cinco años, y era desproporcionadamente gente en edad de trabajar y de tener hijos.

Esto cae sobre una demografía que ya era de las más envejecidas del hemisferio. La
fecundidad está por debajo del reemplazo desde 1978 ---no hay ninguna cohorte en
camino--- y a mediados de la década de 2020 la tasa global de fecundidad rondaba
1,2 a 1,4, de las más bajas del mundo, con más de una cuarta parte de la población
por encima de los sesenta.\unv{} Un sistema de pensiones, un sistema de salud y una
base tributaria dependen todos de una razón entre trabajadores y dependientes, y
esa razón se ha movido más y más rápido en Cuba que en ningún país que no esté en
guerra.

El capítulo~\ref{ch:welfare} trata esto como la cifra más difícil del libro, y el
capítulo~\ref{ch:talent} extrae la conclusión que se sigue: un país que ha perdido
esta cantidad de gente en edad de trabajar no puede diseñarse como si la mano de
obra fuera abundante. Toda propuesta de la Parte III está limitada por la mano de
obra.

\section{Dónde está el país}

En 2026 Cuba es más pobre, más vieja, más pequeña, más oscura y con más hambre que
en ningún momento desde 1994, y por varias medidas que en ningún momento desde
1959.

El producto se ha contraído en la mayoría de los años desde 2019 y la caída
acumulada es grande.\unv{} El azúcar, que fue de 8,5 millones de toneladas en 1970
y de 4 millones tan recientemente como en los años dos mil, produjo zafras medidas
en cientos de miles de toneladas a mediados de la década de 2020 ---menos que en el
siglo XIX--- y Cuba ahora importa azúcar.\unv{} La canasta racionada se entrega
tarde e incompleta, y en 2023 el gobierno pidió al Programa Mundial de Alimentos
asistencia con leche en polvo para niños, cosa que nunca antes había
hecho.\unv{} Las medicinas faltan en las farmacias en proporciones que el propio
ministerio de salud ha reconocido. Los médicos y enfermeros emigran, y los maestros
también. Los apagones continúan.

Frente a eso: las escuelas están abiertas, las clínicas atendidas aunque sea
escasamente, el programa de vacunación sigue funcionando, el crimen violento sigue
siendo bajo para los estándares regionales, y el Estado no ha perdido el control
del territorio. El piso que aguantó en 1993 es más delgado y sigue ahí. Es lo que
la Parte III está escrita para salvar.

\begin{ledger}{la crisis, 2018--2026}
\emph{Logrado.} Una constitución que reconoce la propiedad privada y restablece el
habeas corpus, y un código de familia que legaliza el matrimonio igualitario
aprobado en referendo. Tres vacunas contra la COVID-19 desarrolladas en el país y
una población vacunada con ellas, en solitario en América Latina. Empresas privadas
con personalidad jurídica por fin, tras una década de negativa. Una unificación
monetaria que, cualquiera que fuese su ejecución, puso fin a una ficción contable
de veintiséis años. Capacidad solar iniciada. Y un Estado que ha mantenido las
escuelas abiertas y las clínicas atendidas a través de una contracción que no ha
podido detener.

\emph{Costo.} Una reforma monetaria ejecutada sin reservas, sin respuesta de la
oferta y con el sector privado todavía cerrado, que destruyó los ahorros de la
población por segunda vez en treinta años. Apagones de doce a veinte horas y cuatro
colapsos totales de la red nacional. Producción azucarera por debajo de los niveles
del siglo XIX y una petición al Programa Mundial de Alimentos por leche para niños.
La mayor emigración de la historia cubana ---entre uno y dos millones de personas en
cinco años, de una población de once--- sobre la demografía más envejecida del
hemisferio. Y, el 11 de julio de 2021, condenas de hasta treinta años, menores
incluidos, a cubanos pobres y en su mayoría negros por caminar por la calle diciendo
que no había comida ni electricidad.
\end{ledger}
